%=============================================================================
%  Section 5 -- The mechanism
%  Every number is a macro from macros.tex (authority: paper/NUMBERS.md).
%=============================================================================
\section{The Mechanism: Diffusion Cannot Batch}
\label{sec:mechanism}

Sections~\ref{sec:cost-laws} and~\ref{sec:two-walls} report four separate
phenomena: a stepwise cost law, a length-invariant quality curve, a saturation
knee predicted by its own service time, and a win region bounded in arrival rate
rather than in anything else. This section argues that they are one phenomenon.
The measurement that connects them is the batching multiplier --- how much
aggregate throughput a runtime gains by serving many requests at once instead of
one at a time --- and it is the axis on which the two architectures differ by
nearly two orders of magnitude.

\subsection{The measurement}
\label{sec:mechanism:measurement}

Table~\ref{tab:multipliers} reports, for each model and runtime, throughput at
batch~1, peak throughput over a batch sweep, and their ratio.

\begin{table}[t]
  \centering
  \caption{Batching multipliers. ``Mult.'' is peak throughput divided by
    batch-1 throughput on the same runtime and model. The two diffusion models
    come from different families (\modelDG{} is mixture-of-experts, \modelLLaDA{}
    is dense) and are served by two independent runtimes; both cap near
    $1\times$. Cross-stack ratios are weaker evidence than within-stack ones and
    are labelled as such.}
  \label{tab:multipliers}
  \begin{tabular}{llrrrr}
    \toprule
    Runtime & Model & $L$ & B{=}1 (ev/s) & peak (ev/s) & Mult. \\
    \midrule
    vLLM 0.25.1 & \modelDG{}    & 128 & \needsnum{dgBOneEvsOneTwoEight} & \needsnum{dgPeakEvsOneTwoEight} & \textbf{\dgMultOneTwoEight$\times$} \\
    vLLM 0.25.1 & \modelDG{}    & 512 & \needsnum{dgBOneEvsFiveTwelve}  & \needsnum{dgPeakEvsFiveTwelve}  & \textbf{\dgMultFiveTwelve$\times$} \\
    dInfer 0.1  & \modelLLaDA{} & --- & \multicolumn{2}{c}{flat to batch 64}                            & \textbf{\lladaMultDinfer$\times$} \\
    \midrule
    vLLM 0.25.1 & \modelARthree{} & 128 & \needsnum{arBOneEvsOneTwoEight} & \needsnum{arPeakEvsOneTwoEight} & \textbf{\arMultOneTwoEight$\times$} \\
    vLLM 0.25.1 & \modelARthree{} & 512 & \needsnum{arBOneEvsFiveTwelve}  & \needsnum{arPeakEvsFiveTwelve}  & \textbf{\arMultFiveTwelve$\times$} \\
    \bottomrule
  \end{tabular}
\end{table}

The autoregressive model converts a queue into throughput at
\arMultOneTwoEight$\times$ ($L=128$) and \arMultFiveTwelve$\times$ ($L=512$).
The diffusion models convert a queue into almost nothing:
\dgMultOneTwoEight$\times$ and \dgMultFiveTwelve$\times$ under vLLM,
\lladaMultDinfer$\times$ under dInfer. The batching-capability gap at $L=128$ is
roughly $\batchGap\times$. On the dInfer sweep the shape is as stark as the
ratio: per-batch wall-clock time scaled linearly with batch size across the
entire sweep, so the multiplier is flat out to batch 64 rather than rising and
then plateauing. Adding a sequence adds proportional compute, with nothing
amortised.

Two consequences are worth separating immediately. First, at batch~1 the
diffusion model \emph{leads}: \dgBOneEvs{}~ev/s against \arBOneEvs{}~ev/s. The
autoregressive throughput case is entirely the multiplier, and a multiplier
requires a queue to exist. Where there is no queue --- a single interactive
session, an edge deployment, a per-tenant isolated worker, a serial
investigation chain in which each step must finish before the next begins ---
the multiplier is unavailable by construction and the comparison reverts to the
batch-1 column. Second, the reverse also holds and is the harder half of the
result: wherever a queue does exist, the autoregressive model converts it and
the diffusion model does not, which is exactly the $\lambda$ wall of
Section~\ref{sec:two-walls}.

\subsection{The control that makes this architectural}
\label{sec:mechanism:control}

Taken alone, the \modelDG{} rows would be a bug report. They were measured on a
runtime with a known batching defect (Section~\ref{sec:two-walls:asymmetry}),
and a reviewer would be right to read $\dgMultOneTwoEight\times$ as a property
of \texttt{vllm}'s DiffusionGemma runner rather than of diffusion decoding.

The dInfer row is what changes the claim's status, and it is load-bearing.
\modelLLaDA{} under dInfer 0.1 is a second, independent, \emph{diffusion-native}
runtime: it was written for diffusion language models rather than adapted to
them, it does not share vLLM's scheduler, and \texttt{vllm\#50699} does not
apply to it --- the batch sweep ran to completion with no crash and no
serialization fallback. It reproduces the same cap, \lladaMultDinfer$\times$,
flat to batch 64. The two measurements share neither the runtime, nor the
scheduler, nor the model family --- \modelDG{} is a mixture-of-experts block
diffusion model, \modelLLaDA{} is a dense masked diffusion model --- nor the
underlying vLLM version. What they share is the decoding paradigm, and they
agree on the number that matters. That is what licenses an architectural claim
instead of a defect report, and we would not make the claim without it.

\subsection{What is measured and what is not}
\label{sec:mechanism:honesty}

Three limits on the above, stated before a reader finds them.

\paragraph{The \modelDG{} multiplier is a lower bound, not a ceiling.}
Its batch sweep was capped at batch \dgBatchCap{} by the mixed-phase crash of
Section~\ref{sec:two-walls:asymmetry}: at larger batches the engine dies rather
than returning a number. So $\dgMultOneTwoEight\times$ is a lower bound on what
vLLM could achieve if the runner were fixed, and it is emphatically not a proven
architectural bound. Had the paper only this measurement, the correct conclusion
would be ``vLLM 0.25.1 cannot batch DiffusionGemma,'' which is a much weaker and
much less interesting statement. The dInfer control in
Section~\ref{sec:mechanism:control}, which is uncapped and crash-free, is what
carries the architectural reading.

\paragraph{The dInfer numbers are cache-less.}
dInfer's optimised prefix and vicinity KV-cache paths error on prompts of the
length this task uses (\dinferCachePromptLimit{} tokens and above), so the sweep
ran with caching disabled. This inflates absolute latency substantially and we
do not use dInfer's absolute latencies as evidence anywhere in this section ---
only the ratio between batch sizes, measured within one configuration. We note,
and flag as a prediction rather than a result, that the reported mechanism
implies the multiplier should be insensitive to the cache fix: compute-bound
scaling with batch size is not a caching property.

\paragraph{Sizes are not matched.} \modelDG{} and \modelLLaDA{} are larger than
\modelARthree{}. Multipliers are within-model ratios, so this does not
contaminate the multiplier column; it does contaminate any comparison of
absolute ev/s across rows, including the batch-1 lead noted above, which should
be read as indicative rather than decisive.

\subsection{Proposed mechanism}
\label{sec:mechanism:proposal}

We propose one mechanism for the cap, and we are explicit that it is inferred
from the measurements rather than measured directly.

Continuous batching is profitable for autoregressive decoding because
autoregressive decoding is memory-bandwidth-bound at small batch. Each decode
step touches every weight to produce one token per sequence, and the KV cache
makes all prior context reusable, so the per-step arithmetic is tiny relative to
the bytes moved. Adding sequences to the batch reuses the same weight traffic
across more useful work: the step converts spare arithmetic into throughput at
almost no marginal bandwidth cost. The multiplier is the size of that
conversion.

Bidirectional attention removes both halves of the setup. Every denoising
iteration attends over the whole canvas in both directions, so there is no
prefix whose keys and values remain valid from one iteration to the next in the
way a causal cache requires; the canvas is recomputed rather than extended. The
consequence is that a single denoising step already performs full-sequence
attention and full-sequence feed-forward work for one request. The model is
therefore compute-bound at batch~1 --- there is no bandwidth-bound headroom left
for batching to convert --- and each additional sequence adds proportional
compute. That is precisely the linear per-batch scaling observed on dInfer, and
it predicts a multiplier near $1\times$ on any runtime, which is what both
runtimes show.

\paragraph{The ablation that would confirm it.}
The mechanism above is consistent with the measurements but is not established
by them; the same multiplier would result from several other causes, including
kernel-level inefficiency in both runners. Three experiments would separate
them, and we did not run them.
\begin{enumerate}
  \item \textbf{Roofline instrumentation.} Measure achieved FLOP utilisation and
    HBM bandwidth at batch~1 and at the peak batch for both architectures. The
    proposal predicts the autoregressive model sits near the bandwidth roof at
    batch~1 and moves toward the compute roof as batch grows, while the
    diffusion models sit near the compute roof already at batch~1 and do not
    move. This is the direct test and it is cheap.
  \item \textbf{Causal-mask ablation on identical weights.} Serve the same
    diffusion checkpoint with a causal mask and a KV cache --- decoding the
    canvas block-autoregressively --- and re-measure the multiplier. If it
    recovers, attention bidirectionality is implicated specifically; if it stays
    near $1\times$, the cost lies in the denoising iteration count rather than in
    cache invalidation, and our proposal is wrong in its stated form.
  \item \textbf{Cache-repair re-measurement.} Once dInfer's prefix cache handles
    long prompts, repeat the batch sweep. The proposal predicts the multiplier
    is unchanged. An increase would falsify it.
\end{enumerate}

\subsection{One mechanism, four phenomena}
\label{sec:mechanism:unification}

If the account above is right, everything in
Sections~\ref{sec:cost-laws}--\ref{sec:two-walls} follows from it.

\begin{description}
  \item[The $\sim 1\times$ multiplier.] Direct: there is no bandwidth headroom
    at batch~1 to convert, so a queue buys nothing
    (Section~\ref{sec:mechanism:measurement}).
  \item[Flat latency within a canvas, and hence length-invariance.] The unit of
    work is a fixed number of denoising iterations over a fixed canvas, not a
    per-token step. Cost therefore cannot depend on how many of the canvas's
    positions carry content, which is Eq.~\eqref{eq:law-dg}; and because every
    position is decoded jointly under the same number of iterations regardless
    of $L$, neither can quality, which is the flat over-generation regime of
    Section~\ref{sec:cost-laws:quality}. The same fixed-canvas unit explains the
    truncation failure: a canvas cut short has not converged anywhere, so
    nothing is parseable.
  \item[The $\lambda$ wall.] With no multiplier, server capacity equals
    $1/\text{service time}$ exactly, which is why an M/D/1 model parameterised
    solely by the \dgBase{}\,ms service time predicts the measured knee
    (Eq.~\eqref{eq:mdone}) without a fitted term. A model that could batch would
    not obey that equation.
  \item[Why autoregression wins anything batchable.] Its capacity is set by how
    many requests the scheduler can co-resident, not by how long one takes,
    which is why its lag is flat in $\lambda$ across the whole swept range in
    Figure~\ref{fig:joint-grid}.
\end{description}

The mirror image completes the picture, and it is the reason
Section~\ref{sec:cost-laws:quality} belongs in this paper rather than in an
appendix. Autoregressive cost is $O(\text{tokens emitted})$ --- that is the same
per-token step that makes its batching cheap, seen from the other side. Forcing
it to emit tokens it does not need costs time linearly
(Eq.~\eqref{eq:law-ar}) and, less obviously, costs quality too: the continued
generation conditions and corrupts fields the model has already been asked
merely to copy (\arHostFiveTwelve{} against \dgHostFiveTwelve{} on
\texttt{host} at $L=512$). One architecture pays per token in both currencies;
the other pays per canvas in both. The advantage of a diffusion language model
is not latency. It is invariance to output length, in cost and in quality
alike --- and it is realizable only where batching is unavailable, because
wherever a queue can form, the multiplier of Table~\ref{tab:multipliers} decides
the outcome instead.
